\documentclass[10pt,a4paper]{amsart}
\usepackage[margin=19mm]{geometry}
\usepackage{amsmath,amssymb,mathtools}
\usepackage[scr=rsfs,cal=euler]{mathalfa}
\usepackage{xfrac}
\usepackage[unicode,hidelinks,bookmarks=false]{hyperref}
\newcommand{\ie}{i.e.\ }
\newcommand{\cf}{cf.\ }
\newcommand{\resp}{respectively\ }
\newcommand{\Subsection}{\subsection}
\providecommand{\nobreakdash}{\nobreak\hbox{-}}
\setlength{\emergencystretch}{2em}
\multlinegap0pt
\usepackage{bm}
\usepackage{bbm}
\usepackage{dsfont}
\usepackage{xspace}
\usepackage{cancel}
\usepackage{mathtools}
\usepackage{amsthm}
\makeatletter
\@ifundefined{theorem}{\newtheorem{theorem}{Theorem}}{}
\@ifundefined{proposition}{\newtheorem{proposition}[theorem]{Proposition}}{}
\makeatother
\makeatletter
\expandafter\ifx\csname altproof\endcsname\relax
\newenvironment{altproof}[1]
{\noindent
{\em Proof of {#1}}.}
{\nopagebreak\mbox{}\hfill $\Box$\par\addvspace{0.5cm}}
\fi
\makeatother
\title[Controllability of a Fluid-Structure Interaction System Gove]{Controllability of a Fluid-Structure Interaction System Governed by the Heat and Damped Beam Equations}
\author{Mehdi Badra, J\'er\'emi Dard\'e, Emmanuel Zongo}
\date{Source: arXiv:2506.16880v1 (CC BY 4.0)}
\begin{document}
\maketitle

\noindent\emph{Source and licence.} Controllability of a Fluid-Structure Interaction System Governed by the Heat and Damped Beam Equations. Version arXiv:2506.16880v1 (CC BY 4.0), distributed under the Creative Commons Attribution 4.0 International licence, as recorded in the arXiv licence metadata for this exact version and shown on the abstract page https://arxiv.org/abs/2506.16880v1. Full rights notice: licenses/arxiv-2506.16880-CC-BY-4.0.txt.\par\smallskip

\noindent\textit{Editorial scope.} Counted unit: the Proposition whose source label is \texttt{non-b-op-adj} in the article (source0.tex lines 658--671, source label \texttt{non-b-op-adj}), delivered together with the article's own complete original proof of it (source0.tex lines 672--714, one \texttt{proof} environment of 43 lines). No mathematics is written, repaired or re-derived: statement and proof are the article's own text line for line. Required context retained uncounted: non-b-opDomain (lines 569--571). Every definition, display and label of those lines is retained unchanged. Omitted from this package and not redistributed: 1--568, 572--657, 715--2091. Nothing inside a retained excerpt is omitted or altered: every retained body is delivered exactly as the article's lines are written, so no omission receipt of any kind accompanies this package. Citations: the retained excerpts carry no citation command at all, so no bibliography entry of the article is transcribed and no reference is redistributed. Reference adaptation: none was needed and none was made; every reference inside the delivered unit resolves inside it, so no unresolved reference or citation is produced. Printed slips kept literal: no printed word, symbol, hypothesis or number of the counted statement or of its proof was altered. Layout and class: the article's own class is \texttt{article} with packages including tikz, pdfrender, mathrsfs, amsmath, amsthm, amssymb, amsfonts, amscd, amstext, latexsym, appendix, xspace, xcolor, hyperref; the wrapper is the lane's pinned \texttt{amsart} editorial wrapper at 10pt on A4, which restates the theorem-like environments the retained text needs and adds the article's own macro definitions that the retained text needs (none), with extra wrapper packages (bm, bbm, dsfont, xspace, cancel, mathtools). The delivered numbering is the unit's own. Assets: no retained span contains a figure environment, a graphics-inclusion command, a PDF-page inclusion, a table or a TikZ picture; no image file is copied into this package, the package declares no asset dependency and no image file is redistributed with it. Licence: version arXiv:2506.16880v1 of this article is distributed under the Creative Commons Attribution 4.0 International licence, as recorded in the arXiv licence metadata for this exact version and shown on the abstract page https://arxiv.org/abs/2506.16880v1. A full rights notice is delivered at licenses/arxiv-2506.16880-CC-BY-4.0.txt. Characters: every retained excerpt is pure ASCII, so the delivered engine, XeLaTeX with its default Latin Modern text font, reports no missing character.
\begin{equation}\label{non-b-opDomain}
\mathcal{D}(\mathcal{A})=\left\{(w,\zeta_1,\zeta_2)\in H^2(\Omega)\times H^4(\mathcal{I})\times H^2(\mathcal{I})~:~w=\Lambda \zeta_2~~\text{on}~\partial\Omega\right \}.
\end{equation}

\begin{proposition} 
    The adjoint of the operator $\mathcal{A}$ is given by $\mathcal{D}(\mathcal{A}^*)=\mathcal{D}(\mathcal{A})$ and 
    \begin{equation}\label{non-b-op-adj}
\mathcal{A}^*\begin{bmatrix}
u\\
\eta_1\\
\eta_2
\end{bmatrix}=\begin{bmatrix}
  \Delta u\\
  -\eta_2\\
  A_1\eta_1-A_2\eta_2-\Lambda^*\partial_n u
\end{bmatrix}.
\end{equation}
\end{proposition}

\begin{proof}
First, let us prove $\mathcal{D}(\mathcal{A}^*)\subset \mathcal{D}(\mathcal{A})$. For that, suppose that $[u,\eta_1,\eta_2]\in \mathcal{D}(\mathcal{A}^*)$ and set $[F,G,H]=\mathcal{A}^*[u,\eta_1,\eta_2]\in \mathscr{H}$. We have
$$\left(\mathcal{A}[w,\zeta_1,\zeta_2],[u,\eta_1,\eta_2]\right)_{\mathscr{H}}=\left([w,\zeta_1,\zeta_2], [F,G,H]\right)_{\mathscr{H}}\quad \forall ~[w,\zeta_1,\zeta_2]\in \mathcal{D}(\mathcal{A}).$$
The above equality rewrites
\begin{multline}\label{Eq1proofAstar}
\int_\Omega\Delta w~u+\int_{\mathcal{I}}A_1^{1/2} \zeta_2 ~ A_1^{1/2}\eta_1-\int_{\mathcal{I}}(\Lambda^*\partial_n w+A_1\zeta_1+A_2\zeta_2)\eta_2\\
=\int_\Omega w~F+\int_{\mathcal{I}} A_1^{1/2}\zeta_1~ A_1^{1/2} G + \int_{\mathcal{I}}H~\zeta_2.
\end{multline}

First, with $w=0$ and $\zeta_2=0$ we find
$$\int_{\mathcal{I}}A_1 \zeta_1(G+\eta_2)=0 \quad\forall \zeta_1\in H^4(\mathcal{I}).$$ 
from which we deduce that $\eta_2=-G\in \mathcal{D}(A_1^{1/2})$. 

Next, with $\zeta_1=\zeta_2=0$ and $w\in C^\infty_c(\Omega)$, we obtain
$$\int_\Omega u \Delta w=\int_\Omega F w\quad \forall w\in C^\infty_c(\Omega).$$ 
from which we deduce $\Delta u=F\in L^2(\Omega)$. It implies that the trace and the normal derivative of $u$ are well defined in $H^{-1/2}(\partial\Omega)$ and in $H^{-3/2}(\partial\Omega)$ respectively, and,  that the following formula is true for all $w \in H^2(\Omega)$,
\begin{equation}\label{ipp}
\int_\Omega\Delta w\,u-\int_\Omega \Delta u\,w=\langle\partial_n w,u\rangle_{H^{1/2}(\partial\Omega),H^{-1/2}(\partial\Omega)}-\langle w,\partial_n u\rangle_{H^{3/2}(\partial\Omega),H^{-3/2}(\partial\Omega)}.
\end{equation}
Thus, by setting $F=\Delta u$ in \eqref{Eq1proofAstar}, and with $\zeta_1=0$, $\zeta_2\in H^2(\mathcal{I})$ and $w\in H^2(\Omega)$ such that $w=\Lambda \zeta_2$ on $\partial\Omega$, we deduce
\begin{equation}\label{f1}
\int_\Omega\Delta w\,u-\int_\Omega w\,\Delta u+\int_{\mathcal{I}} A_1^{1/2} \zeta_2 \, A_1^{1/2}\eta_1-\int_{\partial\Omega}\partial_n w\,\Lambda \eta_2-\int_{\mathcal{I}}A_2\zeta_2\,\eta_2=\int_{\mathcal{I}}\zeta_2\,H.
\end{equation}
Then with \eqref{ipp}, the equality \eqref{f1} rewrites:
\begin{multline}\label{f2}
\langle\partial_n w,u\rangle_{H^{1/2}(\partial\Omega),H^{-1/2}(\partial\Omega)}-\langle w,\partial_n u\rangle_{H^{3/2}(\partial\Omega),H^{-3/2}(\partial\Omega)}\\+\int_{\mathcal{I}} A_1^{1/2} \zeta_2 \, A_1^{1/2}\eta_1-\int_{\partial\Omega}\partial_n w\,\Lambda \eta_2-\int_{\mathcal{I}}A_2\zeta_2\,\eta_2=\int_{\mathcal{I}}\zeta_2\,H.
\end{multline}
Let $d\in H^{1/2}(\partial\Omega)$. By using the surjectivity of the map $w\in H^2(\Omega)\mapsto (w,\partial_n w)\in H^{3/2}(\partial\Omega)\times H^{1/2}(\partial\Omega)$ we can choose $w\in H^2(\Omega)$ such that $\partial_n w=d$ on $\partial\Omega$. Hence by applying \eqref{f2} for such a $w$ and with $\zeta_2=0$ on $\mathcal{I}$ we finally deduce that
\begin{equation}\label{f3} 
\langle d,u\rangle_{H^{1/2}(\Gamma_1),H^{-1/2}(\Gamma_1)}-\int_{\partial\Omega}d \,\Lambda \eta_2=0 \quad \forall d\in H^{1/2}(\partial\Omega).
\end{equation}
Then we deduce that $u=\Lambda \eta_2$ on $\partial\Omega$. Moreover, since $\eta_2\in H^2(\mathcal{I})$ it is clear that $u=\Lambda \eta_2\in H^{2}(\partial\Omega)$. Since we also have $\Delta u\in L^2(\Omega)$, elliptic regularity results yield $u\in H^2(\Omega)$. In particular, we have $\partial_n u\in H^{1/2}(\partial\Omega)$ and $\Lambda^*\partial_n u\in H^{1/2}(\mathcal{I})$.

Finally, by coming back to \eqref{f2}, that holds for all $w\in H^2(\Omega)$ and $\zeta_2\in H^2(\mathcal{I})$ such that $w=\Lambda \zeta_2$ on $\partial\Omega$, and by taking into account that $u=\Lambda\eta_2$ on $\partial\Omega$, we find that 
$$
-\int_{\partial\Omega} \Lambda \zeta_2\,\partial_n u+\int_{\mathcal{I}} A_1^{1/2} \zeta_2 \, A_1^{1/2}\eta_1-\int_{\mathcal{I}}A_2\zeta_2\,\eta_2=\int_{\mathcal{I}}\zeta_2\,H
.$$
Then $A_1 \eta_1=H+A_2\eta_2 +\Lambda^*\partial_n u\in L^2(\mathcal{I})$ from which we deduce $\eta_1\in \mathcal{D}(A_1)$. 

If we summarize the above results, we have proved $(u,\eta_1,\eta_2)\in H^2(\Omega)\times \mathcal{D}(A_1)\times \mathcal{D}(A_1^{1/2})$ and $u=\Lambda\eta_2$ on $\partial\Omega$, and that $F=\Delta u$, $G=-\eta_2$ and $H=A_1 \eta_1-A_2\eta_2 -\Lambda^*\partial_n u$. Then with  \eqref{non-b-opDomain} it means that $(u,\eta_1,\eta_2)\in \mathcal{D}(\mathcal{A})$, and we have proved $\mathcal{D}(\mathcal{A}^*)\subset \mathcal{D}(\mathcal{A})$. Moreover, since $(F,G,H)=\mathcal{A^*}(u,\eta_1,\eta_2)$ we have that \eqref{non-b-op-adj} is true. 

To conclude,  the converse inclusion $\mathcal{D}(\mathcal{A})\subset \mathcal{D}(\mathcal{A}^*)$ follows from an application of the Green formula. 
\end{proof}

\end{document}
